\documentclass{article}
\usepackage[utf8]{inputenc}
\usepackage{amsmath, amssymb, amsthm}
\usepackage{booktabs}
\usepackage{hyperref}
\usepackage{natbib}
\usepackage[margin=1in]{geometry}

\newtheorem{proposition}{Proposition}

\title{6998 -- Actionable Project Proposal (Revised)}
\author{Jeremiah Mao, Mateo Juliani}
\date{2026-02-27}

\begin{document}
\maketitle

\textbf{Problem:}
Leaving LLMs static post-deployment leads to performance degradation as the world changes \citep{lazaridou2021mind}. While fine-tuning can update a model's knowledge, it risks overwriting previously learned capabilities \citep{lin2024mitigating}. Several knowledge update strategies have been proposed---including locate-and-edit methods such as AlphaEdit \citep{fang2025alphaedit} and continual alignment methods such as COPR \citep{zhang2024copr}---but they have been evaluated almost exclusively on base or instruction-tuned models \citep{li2024shouldwe}.

In practice, deployed models are often fine-tuned for specific downstream tasks, learning not just facts but \emph{behavioral strategies}. A recent study found that fine-tuning moves activations in directions \emph{nearly orthogonal} to knowledge editing directions, and that AlphaEdit edits are more fragile than MEMIT edits under subsequent fine-tuning \citep{wang2025finetune}. This suggests a fundamental mismatch: weight-space editing methods provide \emph{pointwise} preservation guarantees on a finite set of inputs, but task-tuned models encode learned skills as \emph{distributional strategies} over an unbounded input space.

We investigate whether framing knowledge updates as a \textbf{policy preservation problem}---regularizing the model's output distribution rather than constraining individual weight updates---better preserves task performance in task-tuned models. We study this in the financial domain, where rapidly shifting market dynamics make continual knowledge updates a practical necessity. Consider, for example, a model fine-tuned to decompose financial queries into sub-queries for retrieval-augmented generation (RAG): it must continually absorb new entity and event knowledge while preserving its learned decomposition strategy.

\bigskip

\textbf{Contributions:}
\begin{enumerate}
    \item \textbf{A concrete adaptation of COPR for factual knowledge injection.} COPR \citep{zhang2024copr} was designed for continual preference alignment using human-ranked responses. We adapt it for factual knowledge updates by replacing human preference rankings with \emph{factual correctness rankings}: for each knowledge update, we sample multiple model responses, rank them by ground-truth accuracy, and apply COPR's advantage-weighted policy fitting with task replay regularization. This bridges the continual alignment and knowledge editing literatures.

    \item \textbf{The first systematic comparison of knowledge update strategies on task-tuned models}, including weight-space editing (AlphaEdit), policy-space regularization (COPR-adapted), and SFT baselines with varying degrees of regularization, evaluated across two distinct financial tasks---numerical reasoning (FinQA) and query decomposition for RAG---to test generality.

    \item \textbf{An ablation isolating the contribution of advantage-weighted fitting} beyond simple KL regularization, testing whether COPR's structured policy fitting provides benefits over naive KL-regularized SFT for preserving task performance during knowledge injection.
\end{enumerate}

\bigskip

\textbf{Theoretical Motivation:}

We ground our comparison in a formal distinction between two classes of preservation guarantees.

\textit{Weight-space preservation (AlphaEdit).}
AlphaEdit \citep{fang2025alphaedit} edits a single FFN weight matrix $W \in \mathbb{R}^{d_1 \times d_0}$ by projecting a perturbation $\Delta$ onto the null space of preserved knowledge. Given preserved keys $K_0 \in \mathbb{R}^{d_0 \times n}$ (typically $n = 100{,}000$ triples from Wikipedia), the null-space projector is:
\begin{equation}
    \{U, \Lambda, U^T\} = \mathrm{SVD}(K_0 K_0^T), \quad P = \hat{U}\hat{U}^T
\end{equation}
where $\hat{U}$ are eigenvectors with eigenvalues below a threshold. The sequential editing solution is:
\begin{equation}
    \Delta_{\text{AlphaEdit}} = R \cdot K_1^T \cdot P \cdot (K_p K_p^T P + K_1 K_1^T P + I)^{-1}
\end{equation}
This yields an \textbf{exact but pointwise} guarantee: $(W + \Delta P) K_0 = W K_0$. Outputs are preserved on the finite set $K_0$, but no guarantee extends to unseen inputs. Moreover, the constraint applies to a single FFN layer, while task skills involve coordinated representations across many layers \citep{wang2025finetune}. AlphaEdit has been tested up to ${\sim}3{,}000$ edits; beyond this, general capabilities degrade significantly \citep{li2024shouldwe}.

\textit{Policy-space preservation (COPR, adapted).}
COPR \citep{zhang2024copr} derives the optimal updated policy as:
\begin{equation}
    \pi^*_t(y|x) = \frac{1}{Z_t(x)} \pi_{t-1}(y|x) \cdot \exp\!\left(\frac{r_t(x,y)}{\beta}\right)
\end{equation}
where $\pi_{t-1}$ is the previous policy and $r_t$ encodes the learning signal. Originally, $r_t$ is derived from human preference rankings over candidate responses $\mathcal{Y}^x = \{y_1 \prec \ldots \prec y_J\}$, with a linear advantage $\mathrm{Adv}(x, y_j) = (2j - J_x - 1)/J_x$.

\textbf{Our adaptation:} We replace human preference rankings with \emph{factual correctness rankings}. For each knowledge-update prompt $x$ (e.g., ``Who is the current CEO of Company X?''), we sample $K$ responses from the current model, rank them by ground-truth accuracy (correct answer $\to$ rank $K$; hallucinations $\to$ rank 1), and compute COPR's advantage function over these factual rankings. The training objective combines fitting on new-knowledge responses with regularization on a task replay buffer:
\begin{align}
    \mathcal{L}^{\text{fit}}_t(\theta) &= \mathbb{E}_{x \sim \mathcal{D}_{\text{new}}} \sum_{y \in \mathcal{Y}^x} D_{\mathrm{KL}}\!\left(P_{y,t}(y|x,\theta) \,\|\, P^*_{y,t}(y|x)\right) \\
    \mathcal{L}^{\text{reg}}_t(\theta) &= \mathbb{E}_{x \sim \mathcal{R}_{\text{task}}} D_{\mathrm{KL}}\!\left(P_{y,t}(y|x,\theta) \,\|\, P^*_{\text{task}}(y|x)\right)
\end{align}
where $\mathcal{R}_{\text{task}}$ is a 5\% replay buffer of task training data (FinQA or query decomposition examples, depending on the task model).

This yields a \textbf{distributional} guarantee via Pinsker's inequality. If $\mathbb{E}_{x \sim \mathcal{D}_{\text{task}}}[D_{\mathrm{KL}}(\pi_{\text{task}} \| \pi_\theta)] \leq \epsilon$, then for any bounded task metric $R$ with $\|R\|_\infty \leq 1$:
\begin{equation}
    |\mathbb{E}_{\pi_{\text{task}}}[R] - \mathbb{E}_{\pi_\theta}[R]| \leq \sqrt{2\epsilon}
\end{equation}

\textit{Connection to EWC.}
The Fisher information matrix $F$ is the Hessian of the KL divergence \citep{kirkpatrick2017overcoming}: $D_{\mathrm{KL}}(p_{\theta^*} \| p_\theta) \approx \tfrac{1}{2}(\theta - \theta^*)^T F(\theta^*) (\theta - \theta^*)$. EWC's diagonal approximation is thus a lossy, parameter-level proxy for the output-level KL that COPR controls directly. Recent work shows that KL divergence from the base model predicts catastrophic forgetting with $R^2 = 0.96$ \citep{razak2025rl}, further motivating direct KL control.

\begin{proposition}
Let $\pi_{\text{task}}$ be a task-tuned policy and $\mathcal{D}_{\text{task}}$ the task input distribution.

\textbf{(a) Pointwise (AlphaEdit):} For finite $S = \{k_1, \ldots, k_n\}$: $(W + \Delta P)k_i = Wk_i$ for all $k_i \in S$. No guarantee for $x \notin S$.

\textbf{(b) Distributional (COPR):} If $\mathbb{E}_{x}[D_{\mathrm{KL}}(\pi_{\text{task}} \| \pi_\theta)] \leq \epsilon$, then $|\mathbb{E}_{\pi_{\text{task}}}[R] - \mathbb{E}_{\pi_\theta}[R]| \leq \sqrt{2\epsilon}$ for any bounded $R$.

For task-tuned models---where the learned skill is a distributional strategy---guarantee (b) is the appropriate abstraction.
\end{proposition}

\bigskip

\textbf{Approach:}

\textit{Task models.}
We fine-tune Qwen2.5-3B on two financial tasks that represent distinct behavioral skills:

\begin{enumerate}
    \item \textbf{FinQA} \citep{chen2021finqa}: Multi-step numerical reasoning over financial reports (8,281 QA pairs over SEC filings with tables and text). The model must parse tables, identify quantities, and compose arithmetic programs. Models of comparable scale (phi-3-mini, 3.8B) achieve 73.5\% execution accuracy with LoRA \citep{loukas2024finetuning}. This is an established benchmark that ensures reproducibility.

    \item \textbf{Financial query decomposition for RAG}: Given a financial question (e.g., ``How did the Middle East conflict affect energy sector earnings in Q4 2023?''), the model generates 2--4 sub-queries whose embeddings maximize Recall@10 of target documents in a retrieval corpus. We train this via SFT on decomposition examples, optionally refined with DPO using Recall@10 as the preference signal. This task is more directly sensitive to entity knowledge---generating effective sub-queries requires knowing \emph{which} entities, events, and relationships are relevant---making it a stronger test of whether knowledge updates interact with task performance.
\end{enumerate}

Using two tasks lets us test whether our findings generalize across task types: FinQA tests preservation of a \emph{procedural reasoning} skill, while query decomposition tests preservation of a \emph{strategic generation} skill that depends more directly on entity knowledge.

\textit{Knowledge updates.}
We define updates as structured financial fact triples---entity changes that can be applied by all methods on equal footing:
\begin{itemize}
    \item Leadership changes: \texttt{(``Apple'', ``CEO'', ``Tim Cook'')} $\to$ \texttt{(``Apple'', ``CEO'', ``New Person'')}
    \item Acquisitions: \texttt{(``Company X'', ``acquired\_by'', ``Company Y'')}
    \item Financial metrics: \texttt{(``Company X'', ``Q3\_revenue'', ``\$4.2B'')} $\to$ \texttt{(``\$5.1B'')}
\end{itemize}
We extract these triples from FNSPID \citep{dong2024fnspid} (15.7M timestamped financial news records, 1999--2023) using a temporal split, focusing on entities that also appear in FinQA reports so that representational changes from knowledge injection can plausibly interact with the task. We test at three scales: 200, 1{,}000, and 3{,}000 fact edits (the last being AlphaEdit's tested limit).

\textit{Methods compared.}

\begin{table}[h]
\centering
\small
\begin{tabular}{@{}lll@{}}
\toprule
\textbf{Method} & \textbf{Type} & \textbf{Description} \\
\midrule
No update & Lower bound & Task-tuned model, no knowledge injection \\
Naive SFT & Baseline & SFT on fact QA pairs, no regularization \\
KL-reg SFT & Baseline & SFT $+ \lambda D_{\mathrm{KL}}(\pi_{\text{task}} \| \pi_\theta)$ \\
Mixed replay & Baseline & SFT on new facts + 5\% task data shuffled \\
AlphaEdit & Weight-space & Null-space projected fact editing \\
COPR-adapted & Policy-space & Advantage-weighted fitting + task replay \\
Full retrain & Upper bound & Train from scratch on all data \\
\bottomrule
\end{tabular}
\end{table}

This design isolates three key comparisons:
\begin{itemize}
    \item \textbf{Naive SFT $\to$ KL-reg SFT $\to$ COPR-adapted}: Does KL regularization help? Does advantage-weighted fitting help \emph{beyond} KL regularization?
    \item \textbf{AlphaEdit vs.\ COPR-adapted}: Weight-space vs.\ policy-space, head-to-head on the same fact triples.
    \item \textbf{Scaling}: At what edit count does each method begin degrading task performance?
\end{itemize}

\textit{COPR adaptation pipeline (detailed).}
For each fact triple to inject:
\begin{enumerate}
    \item Convert the triple to a natural-language question (e.g., ``Who is the CEO of Apple?'').
    \item Sample $K = 8$ responses from the current task-tuned model.
    \item Rank responses by factual correctness using the ground-truth answer (exact match $\to$ top rank; partial match $\to$ middle; hallucination $\to$ bottom).
    \item Compute COPR's linear advantage: $\mathrm{Adv}(x, y_j) = (2j - K - 1)/K$.
    \item Compute the renormalized optimal sampling distribution $P^*$.
    \item Train via $\mathcal{L}^{\text{fit}}$ on these ranked responses.
    \item Simultaneously train via $\mathcal{L}^{\text{reg}}$ on a 5\% replay buffer of task training examples (FinQA or query decomposition).
\end{enumerate}

\bigskip

\textbf{Dataset:}
\begin{itemize}
    \item \textbf{Task 1 benchmark:} FinQA \citep{chen2021finqa}---8,281 QA pairs requiring multi-step numerical reasoning over financial reports. Train/dev/test split per the original paper.
    \item \textbf{Task 2 data:} We construct query decomposition training data from FNSPID \citep{dong2024fnspid} by pairing financial questions with gold sub-query decompositions, using the FNSPID news corpus as the retrieval target. We evaluate via Recall@10: given a question, the model's generated sub-queries are embedded and used to retrieve from the corpus, and we measure whether the gold-relevant documents appear in the top 10. Training data is generated using a teacher LLM and filtered by retrieval performance.
    \item \textbf{Knowledge source:} FNSPID---15.7M timestamped financial news records (1999--2023) across 4,775 S\&P 500 companies. We extract structured fact triples from articles after a temporal cutoff, filtered to entities appearing in FinQA reports and/or the retrieval corpus.
    \item \textbf{Fact triple extraction:} We use an LLM to extract (subject, relation, object) triples from post-cutoff FNSPID articles, then verify against ground truth. This is \emph{not} a dataset contribution---it is a preprocessing step to create fair inputs for all methods.
\end{itemize}

\bigskip

\textbf{Metrics:}
For each method, task, and edit scale, we evaluate:
\begin{enumerate}
    \item \textbf{Task preservation:} FinQA execution accuracy (Task 1) and Recall@10 on held-out queries (Task 2).
    \item \textbf{Knowledge absorption:} Accuracy on fact-probe questions about the injected triples (e.g., ``Who is the CEO of Apple?'').
    \item \textbf{Scaling behavior:} Task preservation and knowledge absorption as a function of edit count (200, 1K, 3K).
    \item \textbf{Compute cost:} GPU-hours per update.
\end{enumerate}

We expect Task 2 (query decomposition) to be \emph{more sensitive} to knowledge updates than Task 1 (FinQA), since generating effective sub-queries depends directly on entity knowledge, while numerical reasoning over tables is more procedural. If this differential sensitivity is observed, it strengthens the argument that the choice of update method matters most when the task and knowledge are tightly coupled.

\bigskip

\bibliographystyle{plainnat}
\begin{thebibliography}{99}

\bibitem[Chen et~al.(2021)]{chen2021finqa}
Zhiyu Chen, Wenhu Chen, Charese Smiley, Sameena Shah, Iana Borber, Yanyan Lan, Jiaqi Mu, Joe Leszkowicz, Samuel R.~Bowman, and Wang-Chiew Tan.
\newblock Fin{QA}: A dataset of numerical reasoning over financial data.
\newblock In \emph{EMNLP}, 2021.

\bibitem[Dong et~al.(2024)]{dong2024fnspid}
Zihan Dong, Xinyu Fan, and Zhiyuan Peng.
\newblock {FNSPID}: A comprehensive financial news dataset in time series.
\newblock In \emph{KDD}, 2024.

\bibitem[Fang et~al.(2025)]{fang2025alphaedit}
Junfeng Fang, Houcheng Jiang, Kun Wang, Yunshan Ma, Shi Jie, Xiang Wang, Xiangnan He, and Tat-Seng Chua.
\newblock Alpha{E}dit: Null-space constrained knowledge editing for language models.
\newblock In \emph{ICLR (Outstanding Paper)}, 2025.

\bibitem[Kirkpatrick et~al.(2017)]{kirkpatrick2017overcoming}
James Kirkpatrick, Razvan Pascanu, Neil Rabinowitz, Joel Veness, Guillaume Desjardins, Andrei~A. Rusu, Kieran Milan, John Quan, Tiago Ramalho, Agnieszka Grabska-Barwinska, Demis Hassabis, Claudia Clopath, Dharshan Kumaran, and Raia Hadsell.
\newblock Overcoming catastrophic forgetting in neural networks.
\newblock \emph{PNAS}, 114(13):3521--3526, 2017.

\bibitem[Lazaridou et~al.(2021)]{lazaridou2021mind}
Angeliki Lazaridou, Adhiguna Kuncoro, Elena Gribovskaya, Devang Agrawal, Adam Liska, Tayfun Terzi, Mai Gimenez, Cyprien de~Masson~d'Autume, Tomas Kocisky, Sebastian Ruder, Dani Yogatama, Kris Cao, Susannah Young, and Phil Blunsom.
\newblock Mind the gap: Assessing temporal generalization in neural language models.
\newblock \emph{NeurIPS}, 2021.

\bibitem[Li et~al.(2024)]{li2024shouldwe}
Qi Li, Xiang Yue, and Huan Sun.
\newblock Should we really edit language models? {O}n the evaluation of edited language models.
\newblock In \emph{NeurIPS}, 2024.

\bibitem[Lin et~al.(2024)]{lin2024mitigating}
Yong Lin, Hangyu Lin, Wei Xiong, Shizhe Diao, Jianmeng Liu, Jipeng Zhang, Rui Pan, Haoxiang Wang, Wenbin Hu, Hanning Zhang, Hanze Dong, Renjie Pi, Han Zhao, Nan Jiang, Heng Ji, Yuan Yao, and Tong Zhang.
\newblock Mitigating the alignment tax of {RLHF}.
\newblock \emph{arXiv preprint arXiv:2309.06256}, 2024.

\bibitem[Loukas et~al.(2024)]{loukas2024finetuning}
Gerasimos Loukas, Konstantinos Bougiatiotis, and Anatolia Evangelou.
\newblock Fine-tuning smaller language models for question answering over financial documents.
\newblock \emph{arXiv preprint arXiv:2408.12337}, 2024.

\bibitem[Razak et~al.(2025)]{razak2025rl}
Ahsan Razak, Hao Peng, and Alane Suhr.
\newblock {RL}'s Razor: Why online reinforcement learning forgets less.
\newblock \emph{arXiv preprint arXiv:2509.04259}, 2025.

\bibitem[Wang et~al.(2025)]{wang2025finetune}
Qi Wang, Letian Peng, and Jingbo Shang.
\newblock Can fine-tuning erase your edits?
\newblock \emph{arXiv preprint arXiv:2511.05852}, 2025.

\bibitem[Zhang et~al.(2024a)]{zhang2024copr}
Han Zhang, Lin Gui, Yuanzhao Zhai, Hui Wang, Yu Lei, and Ruifeng Xu.
\newblock {COPR}: Continual learning human preference through optimal policy regularization.
\newblock In \emph{ACL Findings}, 2025.

\bibitem[Zhang et~al.(2024b)]{zhang2024cppo}
Han Zhang, Yu Lei, Lin Gui, Min Yang, Yulan He, Hui Wang, and Ruifeng Xu.
\newblock {CPPO}: Continual learning for reinforcement learning with human feedback.
\newblock In \emph{ICLR}, 2024.

\end{thebibliography}

\end{document}
